\begin{abstract}
Self-evolving agents turn execution trajectories into persistent skills, prompts, or memories that alter future behavior. Existing evaluations usually vary which feedback the updater sees while treating the update mechanism itself as a stable conduit from evidence to state. We show that this assumption can fail. In a controlled SpreadsheetBench skill-evolution substrate, a 48-pair study that replaces winner evidence with a deterministic failed nonwinner yields a small positive mean effect (+2.31 points) but neither superiority nor practical equivalence, with substantial stream heterogeneity. We then localize the instability. A historically strong learned state remains useful when frozen and re-evaluated, but replaying the byte-identical updater evidence does not regenerate the historical effect: two fresh updater realizations yield 0 and $-5.56$ points relative to matched winner-only states. Thus, in this development case, the evidence is not the state: persistent-state synthesis is itself a candidate source of variance. We formalize the resulting causal chain as evidence selection $\rightarrow$ state generation $\rightarrow$ persistent-state semantics $\rightarrow$ downstream behavior, and introduce a variance-controlled typed state compiler that maps learner-visible trajectory evidence to a finite repair representation before canonical compilation. A prospective $2\times2$ bridge crosses evidence source (Winner vs. First-Fail-4) with state generator (native free-form updater vs. constrained compiler), with a score-matched trajectory-blind falsifier and disjoint SCREEN/VALIDATION panels. The paper's central claim is therefore not that failures are inherently better feedback, nor that structured diagnosis is novel, but that state-generation realization is an independently testable bottleneck in persistent agent self-evolution.
\end{abstract}
